\chapter{Equações diferenciais}\label{ch:b2:diffeq}

O primeiro ano resolveu as equações lineares que admitem fórmulas. Este capítulo
fornece o que as fórmulas não conseguem: o \emph{teorema de Cauchy--Lipschitz}
--- existência e unicidade para $y' = f(t, y)$ --- demonstrado com o
teorema do ponto fixo de Banach, exatamente como prometido no
\cref{ch:b2:metric}; e depois a teoria completa dos \emph{sistemas
lineares} $X' = A(t)X + B(t)$, com a \omterm{ex:b2:nvs:matrixexp}{exponencial de matriz} e o
\omterm{def:b2:diffeq:wronskian}{wronskiano} como motores de cálculo.

\section{O teorema de Cauchy--Lipschitz}

\begin{theorem}[Cauchy--Lipschitz, versão globalmente lipschitziana]\label{thm:b2:diffeq:cauchylipschitz}
Seja $I$ um segmento e $f \colon I \times \R^n \to \R^n$ \omterm{def:b2:metric:continuity}{contínua}
e \emph{\omterm{def:b2:metric:continuity}{lipschitziana} na segunda variável}, \omterm{def:b2:funcseq:def}{uniformemente} na
primeira: $\norm{f(t, y) - f(t, z)} \leq k\,\norm{y - z}$ para todos $t
\in I$. Então, para todo $(t_0, y_0) \in I \times \R^n$, o problema
de Cauchy\index{teorema de Cauchy--Lipschitz}
\[
y' = f(t, y), \qquad y(t_0) = y_0
\]
tem exatamente uma solução $y \colon I \to \R^n$ de classe $C^1$.
\end{theorem}

\begin{proof}
\emph{Reformulação.} Uma $y$ \omterm{def:b2:metric:continuity}{contínua} resolve o problema se e somente se
satisfaz a equação integral
\[
y(t) = y_0 + \int_{t_0}^{t} f\bigl(s, y(s)\bigr)\,\dd s
=: T(y)(t)
\]
(teorema fundamental do cálculo nos dois sentidos; uma solução \omterm{def:b2:metric:continuity}{contínua}
da equação integral é automaticamente $C^1$).

\emph{Uma contração, após renormar.} No \omterm{def:b2:nvs:banach}{espaço de Banach} $E =
C(I, \R^n)$ com a \omterm{def:b2:nvs:norm}{norma} \emph{ponderada}
\[
N(y) = \sup_{t \in I}\; \eu^{-2k\abs{t - t_0}}\,\norm{y(t)} ,
\]
(equivalente à \omterm{def:b2:nvs:norm}{norma} do sup: o peso é majorado e minorado
no segmento $I$, de modo que $E$ continua \omterm{def:b2:metric:complete}{completo}), estime para $y, z \in
E$ e, digamos, $t \geq t_0$:
\[
\norm{T(y)(t) - T(z)(t)}
\leq \int_{t_0}^{t} k\,\norm{y(s) - z(s)}\,\dd s
\leq k\,N(y - z)\int_{t_0}^{t} \eu^{2k(s - t_0)}\dd s
\leq \frac{N(y-z)}{2}\,\eu^{2k(t - t_0)} .
\]
Multiplicando por $\eu^{-2k(t - t_0)}$ e tomando o sup (o caso $t
< t_0$ é \omterm{def:b2:quadratic:adjoint}{simétrico}): $N\bigl(T(y) - T(z)\bigr) \leq \frac12 N(y -
z)$: $T$ é uma $\frac12$-contração do \omterm{def:b2:metric:complete}{completo} $(E, N)$. O
teorema do ponto fixo de Banach (\cref{thm:b2:metric:banach}) fornece um
único ponto fixo: a única solução.
\end{proof}

\begin{remark}
Para $f$ apenas $C^1$ (localmente \omterm{def:b2:metric:continuity}{lipschitziana}), o teorema vale
\emph{localmente}, com solução maximal num intervalo \omterm{def:b2:metric:topology}{aberto} maximal;
as soluções podem explodir em tempo finito ($y' = y^2$, $y(0) = 1$: $y(t)
= \frac{1}{1-t}$, desaparecida em $t = 1$). A hipótese globalmente \omterm{def:b2:metric:continuity}{lipschitziana}
é o que compra o segmento inteiro. Duas consequências que vale a pena gravar:
as curvas solução de uma EDO com campo \omterm{def:b2:metric:continuity}{lipschitziano} \emph{nunca se cruzam};
e a função nula é a única solução que se anula em algum ponto de uma
equação linear homogênea.
\end{remark}

\begin{example}[A unicidade é um teorema: um campo que vaza]\label{ex:b2:diffeq:leaky}
Considere $y' = 2\sqrt{\abs y}$ com $y(0) = 0$. A função nula
a resolve; e também
\[
y(t) = \begin{cases} 0 & t \leq 0,\\ t^2 & t \geq 0,
\end{cases}
\]
que é $C^1$ (as duas peças têm derivada $0$ no ponto de
colagem) e satisfaz $y'(t) = 2t = 2\sqrt{t^2}$ para $t > 0$ ---
de fato, atrasar a decolagem dá uma solução para \emph{cada}
instante de largada $c \geq 0$: infinitas soluções pelos
mesmos dados iniciais. Nenhuma contradição com o
\cref{thm:b2:diffeq:cauchylipschitz}: perto de $y = 0$,
\[
\frac{\abs{2\sqrt y - 2\sqrt z}}{\abs{y - z}}
= \frac{2}{\sqrt y + \sqrt z} \longrightarrow +\infty ,
\]
o campo não é \omterm{def:b2:metric:continuity}{lipschitziano} em $y$, e o teorema se cala.
Lição final: a leitura física é um balde esvaziando
sob gravidade, rodado ao contrário --- a partir do estado vazio, não
se consegue dizer quando ele começou a encher; o determinismo das EDOs é
exatamente a condição \omterm{def:b2:metric:continuity}{lipschitziana}, não uma lei da natureza.
\end{example}

\section{Sistemas lineares}

\begin{theorem}[Estrutura dos sistemas lineares]\label{thm:b2:diffeq:linear}
Sejam $A \colon I \to \mathcal{M}_n(\R)$ e $B \colon I \to \R^n$
\omterm{def:b2:metric:continuity}{contínuas} num intervalo $I$. Para todo $(t_0, X_0)$ o problema
\[
X' = A(t)X + B(t), \qquad X(t_0) = X_0
\]
tem exatamente uma solução em \emph{todo} o $I$. As soluções do
sistema homogêneo ($B = 0$) formam um espaço vetorial $\mathcal{S}_H$ de
dimensão exatamente $n$, e a avaliação $X \mapsto X(t_0)$ é um
isomorfismo $\mathcal{S}_H \to \R^n$; solução geral $=$
particular $+$ homogênea.
\end{theorem}

\begin{proof}
Em todo segmento $J \subseteq I$ que contenha $t_0$: $f(t, X) =
A(t)X + B(t)$ é \omterm{def:b2:metric:continuity}{contínua} e \omterm{def:b2:metric:continuity}{lipschitziana} em $X$ com constante $k
= \sup_J \vertiii{A(t)}$ (finita: \omterm{def:b2:metric:continuity}{contínua} num segmento):
o~\cref{thm:b2:diffeq:cauchylipschitz} se aplica em $J$; fazendo $J$
esgotar $I$, a unicidade cola as soluções numa só em $I$.
A linearidade do conjunto solução e da aplicação avaliação é clara;
a avaliação é bijetiva por existência (sobrejetiva) e unicidade
(injetiva): $\dim \mathcal{S}_H = n$. A estrutura afim é o argumento do
primeiro ano, palavra por palavra.
\end{proof}

\begin{example}[O isomorfismo de avaliação, concretamente]\label{ex:b2:diffeq:evaliso}
Para $y'' + y = 0$, visto como o sistema $X' =
\begin{pmatrix} 0 & 1\\ -1 & 0\end{pmatrix}X$ com $X = (y,
y')$: o teorema diz que o espaço de soluções é um plano e que
$X \mapsto X(0) = (y(0), y'(0))$ é um isomorfismo sobre
$\R^2$. As soluções $\cos$ e $\sin$ avaliam em $(1, 0)$
e $(0, 1)$ --- a base canônica de $\R^2$ --- de modo que elas formam
uma base do espaço de soluções, e \emph{toda} solução é
\[
y(t) = y(0)\cos t + y'(0)\sin t ,
\]
com os coeficientes lidos diretamente nos dados iniciais, sem
sistema linear a resolver. Lição final: escolher o
\omterm{def:b2:diffeq:wronskian}{sistema fundamental} cujos valores iniciais são a base
canônica (aqui $\cos, \sin$) é exatamente escolher as colunas de
$\eu^{tA}$; o isomorfismo de avaliação é a razão pela qual as condições
iniciais parametrizam as trajetórias --- o conteúdo geométrico
de ``dinâmica determinística'' para equações lineares.
\end{example}

\begin{definition}[Wronskiano]\label{def:b2:diffeq:wronskian}
Para soluções $X_1, \dots, X_n$ do sistema homogêneo, o
\emph{wronskiano}\index{wronskiano} é $W(t) = \det\bigl(X_1(t),
\dots, X_n(t)\bigr)$. Pelo isomorfismo acima, ou $W$ se anula
identicamente (a família é ligada) ou nunca se anula (um \emph{sistema
fundamental}\index{sistema fundamental}); quantitativamente, $W' = \operatorname{tr}\bigl(A(t)\bigr)
W$, de modo que
\[
W(t) = W(t_0)\,\exp\Bigl(\int_{t_0}^{t}
\operatorname{tr} A(s)\,\dd s\Bigr)
\quad \text{(fórmula de Liouville)}.
\]
\end{definition}

\begin{example}[Liouville conferido numa equação de Euler]\label{ex:b2:diffeq:eulerwronskian}
Em $\intoo{0}{\infty}$, a equação $t^2y'' + ty' - y = 0$ tem
as soluções $y_1(t) = t$ e $y_2(t) = \frac1t$ (substitua).
Seu \omterm{def:b2:diffeq:wronskian}{wronskiano}:
\[
W(t) = \det\begin{pmatrix} t & \tfrac1t\\[2pt]
1 & -\tfrac{1}{t^2}\end{pmatrix}
= -\frac1t - \frac1t = -\frac2t ,
\]
nunca nulo: um \omterm{def:b2:diffeq:wronskian}{sistema fundamental}. Confira agora Liouville: na
forma normalizada $y'' + \frac1t\,y' - \frac{1}{t^2}\,y = 0$, a
matriz companheira $A(t) = \begin{pmatrix} 0 & 1\\ \frac{1}{t^2}
& -\frac1t\end{pmatrix}$ tem traço $-\frac1t$, logo
\[
W(t) = W(1)\exp\Bigl(-\int_1^t\frac{\dd s}{s}\Bigr)
= -2\,\eu^{-\ln t} = -\frac2t . \checkmark
\]
Lição final: Liouville prevê a \emph{forma} do
\omterm{def:b2:diffeq:wronskian}{wronskiano} antes de qualquer solução ser conhecida --- aqui, que $W$ tem de
ser $\frac{c}{t}$; é isso que move o método de
redução de \omterm{def:b2:structures:generated}{ordem}
(\cref{prop:b2:diffeq:secondorder}), em que conhecer $y_1$ e
a forma do \omterm{def:b2:diffeq:wronskian}{wronskiano} determina $y_2$ por uma quadratura.
\end{example}

\begin{proof}[Demonstração da fórmula de Liouville]
$W(t) = \det M(t)$ com $M' = AM$. Derivando o \omterm{def:b2:linalg:det}{determinante}
como função multilinear das colunas,
\[
W'(t) = \sum_j \det(X_1, \dots, X_j', \dots, X_n)
= \sum_j \det(X_1, \dots, AX_j, \dots, X_n) .
\]
Ora, a aplicação $(C_1, \dots, C_n) \mapsto \sum_j \det(C_1, \dots,
AC_j, \dots, C_n)$ é $n$-linear e alternada (com duas colunas
iguais $C_i = C_k$, os termos $j \notin \{i, k\}$ se anulam de saída,
e os termos $j = i$ e $j = k$ se cancelam aos pares após uma troca de
colunas): pelo teorema de unicidade
(\cref{thm:b2:linalg:detspace}) ela vale $c \cdot \det$, com $c$ lido
nas colunas canônicas: $c = \sum_j \det(e_1, \dots, Ae_j, \dots,
e_n) = \sum_j a_{jj} = \operatorname{tr} A$. Logo $W' =
\operatorname{tr}\bigl(A(t)\bigr)W$: uma EDO linear escalar, resolvida pela
fórmula do primeiro ano.
\end{proof}

\section{Coeficientes constantes: a exponencial de matriz}

\begin{theorem}\label{thm:b2:diffeq:matrixexp}
Para $A \in \mathcal{M}_n(\R)$ (ou $\C$), a exponencial $\eu^{tA}
= \sum_k \frac{(tA)^k}{k!}$ (\cref{ex:b2:nvs:matrixexp}) satisfaz:
$t \mapsto \eu^{tA}$ é $C^1$ (na verdade $C^\infty$)
com\index{exponencial de matriz}
\[
\frac{\dd}{\dd t}\,\eu^{tA} = A\,\eu^{tA} = \eu^{tA}A ,
\qquad
\eu^{(s+t)A} = \eu^{sA}\,\eu^{tA},
\qquad
(\eu^{A})^{-1} = \eu^{-A} ;
\]
e $\eu^{A + B} = \eu^A\eu^B$ \emph{quando} $AB = BA$. O problema
de Cauchy $X' = AX$, $X(0) = X_0$ tem a única solução $X(t) =
\eu^{tA}X_0$; com um termo fonte, vale a fórmula de \emph{variação
das constantes}:
\[
X(t) = \eu^{(t - t_0)A}X_0 + \int_{t_0}^{t} \eu^{(t-s)A}B(s)\,\dd
s .
\]
\end{theorem}

\begin{proof}
\emph{\omterm{def:b2:diffcalc:differential}{Diferenciabilidade}:} a série $\sum \frac{t^kA^k}{k!}$ e
sua série derivada termo a termo $\sum \frac{t^{k-1}A^k}{(k-1)!} = A\sum
\frac{(tA)^{k-1}}{(k-1)!}$ convergem \omterm{def:b2:funcseq:series}{normalmente} em todo segmento
(\omterm{def:b2:nvs:norm}{normas} $\leq \frac{(\abs t\,\vertiii A)^k}{k!}$): derive
termo a termo (\cref{thm:b2:funcseq:seriestransfer}, com valores vetoriais).
As duas \omterm{def:b2:structures:generated}{ordens} $A\eu^{tA}$ e $\eu^{tA}A$ coincidem, pois toda
soma parcial comuta com $A$.

\emph{Lei de grupo:} para $A, B$ que comutam, o \omterm{thm:b2:series:fubini}{produto de Cauchy} das
duas séries exponenciais se reorganiza pelo binômio de Newton exatamente
como no~\cref{ex:b2:series:exp} (a convergência absoluta na \omterm{def:b2:structures:algebra}{álgebra} de
Banach o justifica): $\eu^{A+B} = \eu^A\eu^B$; com $B = sA$ isso
dá a lei de grupo a um parâmetro, e $B = -A$ a inversa.

\emph{Problema de Cauchy:} $X(t) = \eu^{tA}X_0$ o resolve (derive);
a unicidade pelo~\cref{thm:b2:diffeq:linear}. Variação das constantes:
ponha $Y(t) = \eu^{-tA}X(t)$; derivando, $Y' = \eu^{-tA}(X' -
AX) = \eu^{-tA}B(t)$; integre de $t_0$ a $t$ e multiplique
de volta por $\eu^{tA}$.
\end{proof}

\begin{method}[Calcular $\eu^{tA}$]\label{met:b2:diffeq:computeexp}
Reduza $A$ (\cref{ch:b2:reduction}): se $A = PDP^{-1}$ é diagonal,
$\eu^{tA} = P\,\eu^{tD}P^{-1}$ com $\eu^{tD}$ diagonal de
$\eu^{t\lambda_i}$; em geral use Dunford $A = D + N$
(que comutam): $\eu^{tA} = \eu^{tD}\,\eu^{tN}$ com $\eu^{tN}$ um
\emph{polinômio} em $t$ (a nilpotência trunca a série). Os
\omterm{def:b2:reduction:eigen}{autovalores} complexos se emparelham em blocos rotação vezes exponencial
(\cref{exo:b2:diffeq:5}).
\end{method}

\begin{remark}[Armadilhas comuns]
\emph{(i) $\eu^{A+B} \neq \eu^A\eu^B$ sem comutação:}
tome $A = \begin{pmatrix} 0 & 1\\ 0 & 0\end{pmatrix}$, $B =
\begin{pmatrix} 0 & 0\\ 1 & 0\end{pmatrix}$. Então $\eu^A = I +
A$, $\eu^B = I + B$ (nilpotência), logo
\[
\eu^A\eu^B = \begin{pmatrix} 2 & 1\\ 1 & 1\end{pmatrix},
\qquad\text{enquanto}\qquad
\eu^{A+B} = \cosh(1)\,I + \sinh(1)\,(A + B)
= \begin{pmatrix} \cosh 1 & \sinh 1\\ \sinh 1 & \cosh 1
\end{pmatrix},
\]
usando $(A+B)^2 = I$; e $\cosh 1 \approx 1.54 \neq 2$. A
lei de grupo do~\cref{thm:b2:diffeq:matrixexp} carrega uma hipótese
de verdade. \emph{(ii) Intuição não linear em terreno linear:}
as soluções de um sistema \emph{linear} com coeficientes \omterm{def:b2:metric:continuity}{contínuos}
vivem no intervalo inteiro
(\cref{thm:b2:diffeq:linear}) --- se uma candidata a solução
explode dentro de $I$, ou a equação não era linear ou o
cálculo está errado; reciprocamente, para equações não lineares
nunca prometa globalidade sem argumento ($y' = y^2$).
\emph{(iii) Dividir pela incógnita:} separar variáveis em
$y' = y(1-y)$ descarta silenciosamente as soluções constantes $0$ e
$1$ --- exatamente as que organizam a reta de fase
(\cref{exo:b2:diffeq:3}); liste primeiro as soluções constantes.
\emph{(iv) Dados iniciais fixam vetores, não escalares:} uma equação
escalar de ordem $n$ precisa de $n$ condições ($y, y', \dots$ em
$t_0$); casar apenas $y(t_0)$ deixa uma família a
$(n-1)$ parâmetros, fonte clássica de constantes
``perdidas''.
\end{remark}

\begin{example}[Uma exponencial $3\times3$ por Dunford]\label{ex:b2:diffeq:dunford3}
Resolva $X' = AX$ para $A = \begin{pmatrix} 2 & 1 & 0\\ 0 & 2 &
0\\ 0 & 0 & 3\end{pmatrix}$. Dunford por blocos: $A = D + N$
com $D = \operatorname{diag}(2, 2, 3)$ e $N = E_{12}$, que
comutam ($N$ vive dentro do bloco de \omterm{def:b2:reduction:eigen}{autovalor} $2$) e $N^2 =
0$:
\[
\eu^{tA} = \eu^{tD}\,\eu^{tN}
= \begin{pmatrix}
\eu^{2t} & t\,\eu^{2t} & 0\\
0 & \eu^{2t} & 0\\
0 & 0 & \eu^{3t}
\end{pmatrix} .
\]
A solução geral se lê coluna a coluna: $X(t) =
\bigl(\eu^{2t}(x_0 + ty_0),\ \eu^{2t}y_0,\
\eu^{3t}z_0\bigr)$. Verificações de bom senso: em $t = 0$ a matriz é
$I$; seu \omterm{def:b2:linalg:det}{determinante} vale $\eu^{7t} =
\eu^{t\operatorname{tr}A}$, como Liouville exige; e o
fator $t$ aparece exatamente onde o \omterm{def:b2:reduction:eigen}{autovalor} $2$ é
defeituoso. Lição final: polinômios vezes exponenciais não são
um chute a decorar --- são as séries truncadas
$\eu^{tN}$, e seu grau é limitado pelo índice de
nilpotência, nunca mais.
\end{example}

\begin{method}[Resolver $X' = AX + B(t)$, do começo ao fim]\label{met:b2:diffeq:solvelinear}
\begin{enumerate}
  \item \omterm{def:b2:reduction:eigen}{Espectro} de $A$; depois $\eu^{tA}$ pelo
        \cref{met:b2:diffeq:computeexp} (diagonalize; ou
        Dunford, como no~\cref{ex:b2:diffeq:dunford3}; ou um
        truque polinomial como o do $A^2 = -I$).
  \item Uma solução particular: a variação das constantes
        $\int_{t_0}^t\eu^{(t-s)A}B(s)\dd s$ sempre funciona;
        para $B$ do tipo exponencial vezes polinômio, um palpite da mesma
        forma (com grau elevado em caso de \omterm{pb:b2:diffeq:1}{ressonância},
        \cref{exo:b2:diffeq:10}) é mais rápido.
  \item Solução geral $= \eu^{(t-t_0)A}X_0 +$ particular;
        case os dados iniciais \emph{por último}, na fórmula
        \omterm{def:b2:metric:complete}{completa}.
  \item Verificações de bom senso: $X(t_0)$ correto; o crescimento da parte homogênea
        casa com as partes reais dos \omterm{def:b2:reduction:eigen}{autovalores}
        (\cref{exo:b2:diffeq:8}); e o $\det$ de uma matriz
        fundamental obedece a Liouville.
\end{enumerate}
\end{method}

\begin{example}[Um retrato de fase]\label{ex:b2:diffeq:phaseportrait}
$X' = AX$ com $A = \begin{pmatrix} 0 & 1\\ -1 & 0\end{pmatrix}$:
$A^2 = -I$, de modo que a série se separa em
\[
\eu^{tA} = (\cos t)\,I + (\sin t)\,A
= \begin{pmatrix} \cos t & \sin t\\ -\sin t & \cos t
\end{pmatrix} :
\]
as trajetórias são círculos percorridos no sentido horário --- o oscilador
harmônico $x'' + x = 0$ em roupagem de primeira ordem. \omterm{def:b2:reduction:eigen}{Autovalores} $\pm\iu$ no
eixo imaginário: um \emph{centro}. Mais geralmente, as partes reais dos
\omterm{def:b2:reduction:eigen}{autovalores} de $A$ decidem o crescimento ou o decaimento de $\norm{X(t)}$
(\cref{exo:b2:diffeq:8}).
\end{example}

\begin{omfigure}
\begin{tikzpicture}
\begin{axis}[omaxis, width=6.4cm, height=6.4cm,
    xmin=-2.3, xmax=2.3, ymin=-2.3, ymax=2.3,
    xtick={-1,1}, ytick={-1,1}]
  \addplot[omDef, thick, domain=0:360, samples=90, smooth]
    ({cos(x)}, {-sin(x)});
  \addplot[omDef!60, thick, domain=0:360, samples=90, smooth]
    ({1.7*cos(x)}, {-1.7*sin(x)});
  \draw[-Stealth, omThm, thick] (axis cs:1,0) -- (axis cs:0.995,-0.12);
  \draw[-Stealth, omThm, thick] (axis cs:1.7,0) -- (axis cs:1.695,-0.12);
\end{axis}
\end{tikzpicture}
\hspace{6mm}
\begin{tikzpicture}
\begin{axis}[omaxis, width=6.4cm, height=6.4cm,
    xmin=-2.3, xmax=2.3, ymin=-2.3, ymax=2.3,
    xtick={-1,1}, ytick={-1,1}]
  % A = diag(-1,-2)-type stable node: trajectories y = c x^2
  \addplot[omMeth, thick, domain=-1.5:1.5, samples=60] {x^2};
  \addplot[omMeth, thick, domain=-1.5:1.5, samples=60] {-x^2};
  \addplot[omMeth, thick, domain=-2.2:2.2, samples=2] {0};
  \addplot[omMeth!60, thick, domain=-1.2:1.2, samples=60] {2*x^2};
  \draw[-Stealth, omThm, thick] (axis cs:1.2,1.44) --
    (axis cs:1.05,1.1);
  \draw[-Stealth, omThm, thick] (axis cs:-1.2,-1.44) --
    (axis cs:-1.05,-1.1);
\end{axis}
\end{tikzpicture}

{\small Dois retratos de fase lineares. À esquerda: um \emph{centro}
(\omterm{def:b2:reduction:eigen}{autovalores} $\pm\iu$) --- órbitas circulares fechadas do oscilador
harmônico. À direita: um \emph{nó estável} (\omterm{def:b2:reduction:eigen}{autovalores} $-1, -2$) ---
todas as trajetórias caem na origem tangencialmente à direção própria
lenta.}
\end{omfigure}

\begin{omfigure}
\begin{tikzpicture}
\begin{axis}[omaxis, width=9.6cm, height=6.4cm,
    xmin=-3.2, xmax=3.2, ymin=-1.6, ymax=2.6,
    xtick={0}, ytick={0},
    xlabel={$\tau = \operatorname{tr}A$},
    ylabel={$\delta = \det A$}]
  \addplot[omThm, very thick, domain=-3:3, samples=100]
    {x^2/4};
  \node[font=\small, omThm] at (axis cs:2.45,2.1)
    {$\delta = \tau^2/4$};
  \node[font=\small] at (axis cs:0,-0.9) {selas};
  \node[font=\small] at (axis cs:-2.45,0.7)
    {nós estáveis};
  \node[font=\small] at (axis cs:2.45,0.7)
    {nós instáveis};
  \node[font=\small] at (axis cs:-1.05,1.9)
    {espirais estáveis};
  \node[font=\small] at (axis cs:1.05,1.9)
    {espirais instáveis};
  \node[font=\small, omDef] at (axis cs:0,2.35) {centros};
  \draw[omDef, very thick] (axis cs:0,0) -- (axis cs:0,2.55);
\end{axis}
\end{tikzpicture}

{\small O plano traço--determinante para $X' = AX$ em dimensão
$2$: abaixo do eixo horizontal, selas; entre o eixo e
a parábola $\delta = \tau^2/4$, nós; dentro da parábola,
espirais; no semieixo positivo de $\delta$, centros. O problema
de fim de semana demonstra essa classificação e segue uma reta
vertical dela --- o oscilador amortecido --- até a \omterm{pb:b2:diffeq:1}{ressonância}.}
\end{omfigure}

\begin{remark}[Onde isso é usado]
Os sistemas lineares são o modelo local de tudo o que é não linear:
perto de um equilíbrio, um campo de vetores suave se comporta (nos
casos hiperbólicos) como sua linearização, cujo retrato o
plano traço--determinante classifica. O problema de fim de semana trabalha
a história do oscilador por inteiro --- amortecimento, forçamento,
\omterm{pb:b2:diffeq:1}{ressonância} e os teoremas de comparação de Sturm para coeficientes
variáveis --- a matemática por trás de amortecedores, circuitos
de corrente alternada e lacunas espectrais. O volume do terceiro ano de graduação volta
com a teoria qualitativa (fluxos, estabilidade, integrais
primeiras) em variedades.
\end{remark}

\section{Segunda ordem com coeficientes variáveis}

\begin{proposition}\label{prop:b2:diffeq:secondorder}
A equação $y'' + a(t)y' + b(t)y = c(t)$ ($a, b, c$ \omterm{def:b2:metric:continuity}{contínuas} em
$I$) é o sistema $X' = A(t)X + B(t)$ para $X = (y, y')$: as soluções
existem e são únicas em todo o $I$ para quaisquer dados iniciais $(y(t_0),
y'(t_0))$; as soluções homogêneas formam um plano. Se uma solução
homogênea $y_1$ que não se anula é conhecida, uma segunda independente
acha-se por \emph{redução de \omterm{def:b2:structures:generated}{ordem}}: pôr $y = y_1 z$ transforma a
equação homogênea numa equação de primeira ordem para $z'$, resolvida
por quadraturas.
\end{proposition}

\begin{proof}
A forma de sistema e o~\cref{thm:b2:diffeq:linear} dão tudo o que é
estrutural. Redução: substituindo $y = y_1z$,
\[
y_1 z'' + (2y_1' + a y_1)z' + \underbrace{(y_1'' + ay_1' +
by_1)}_{=\,0}\,z = 0 :
\]
uma equação linear de primeira ordem em $u = z'$, resolúvel pela fórmula do
primeiro ano; integrar $u$ dá $z$, logo $y_2 = y_1 z$,
independente de $y_1$ sempre que $z$ não é constante.
\end{proof}

\begin{example}\label{ex:b2:diffeq:lowering}
$t^2y'' - 2y = 0$ em $\intoo{0}{\infty}$: $y_1 = t^2$ é uma
solução. Substitua $y = t^2z$: de $y' = t^2z' + 2tz$ e $y''
= t^2z'' + 4tz' + 2z$,
\[
t^2y'' - 2y = t^4 z'' + 4t^3z' = 0,
\qquad\text{i.e.}\qquad \frac{z''}{z'} = -\frac4t :
\]
$z' = t^{-4}$ (a menos de constante), $z = -\frac{1}{3t^3}$ e $y_2
= t^2z = -\frac{1}{3t}$. Solução geral: $y = \alpha t^2 +
\frac{\beta}{t}$.
\end{example}

\section{Exercícios}

\begin{exercise}[$\star$]\label{exo:b2:diffeq:1}
Resolva $X' = AX$, $X(0) = (1, 0)^{\mathsf T}$, para $A =
\begin{pmatrix} 1 & 1\\ 0 & 2\end{pmatrix}$ (diagonalize) e $A =
\begin{pmatrix} 2 & 1\\ 0 & 2 \end{pmatrix}$ (Dunford).
\end{exercise}

\begin{exercise}[$\star$]\label{exo:b2:diffeq:2}
Quais problemas de Cauchy têm solução global única em $\R$ pelo
\cref{thm:b2:diffeq:cauchylipschitz}? $y' = \sin(ty)$; $\;y' =
y^2$; $\;y' = \abs y$. Para o último, resolva explicitamente com $y(0) =
0$ e $y(0) = 1$.
\end{exercise}

\begin{exercise}[$\star$]\label{exo:b2:diffeq:3}
Prove que duas soluções maximais distintas de $y' = f(t,y)$ ($f$
\omterm{def:b2:metric:continuity}{lipschitziana} em $y$) nunca assumem o mesmo valor no mesmo instante, e
deduza que as soluções de $y' = y(1 - y)$ que começam em
$\intoo{0}{1}$ permanecem em $\intoo{0}{1}$ para sempre.
\end{exercise}

\begin{exercise}[$\star\star$]\label{exo:b2:diffeq:4}
Calcule $\eu^{tA}$ para $A = \begin{pmatrix} 3 & 1\\ -1 &
1\end{pmatrix}$ \emph{(Dunford: $(A - 2I)^2 = 0$)} e resolva $X' =
AX + \begin{pmatrix} \eu^{2t}\\ 0\end{pmatrix}$, $X(0) = 0$, por
variação das constantes.
\end{exercise}

\begin{exercise}[$\star\star$]\label{exo:b2:diffeq:5}
Para $A = \begin{pmatrix} \alpha & -\beta\\ \beta &
\alpha\end{pmatrix}$, prove que $\eu^{tA} =
\eu^{\alpha t}\begin{pmatrix} \cos\beta t & -\sin\beta t\\
\sin\beta t & \cos\beta t\end{pmatrix}$ --- trajetórias em espiral ---
de duas maneiras: pela série (escreva $A = \alpha I + \beta J$, $J^2 =
-I$) e pela identificação complexa $z' = (\alpha +
\iu\beta)z$.
\end{exercise}

\begin{exercise}[$\star\star$]\label{exo:b2:diffeq:6}
(Lema de Gronwall) Seja $u$ \omterm{def:b2:metric:continuity}{contínua} não negativa com $u(t)
\leq C + k\int_{t_0}^{t} u(s)\,\dd s$ em $\intco{t_0}{T}$. Prove que
$u(t) \leq C\,\eu^{k(t - t_0)}$ \emph{(derive $v(t) =
\eu^{-kt}\int_{t_0}^t u$)}. Deduza de novo a unicidade em
Cauchy--Lipschitz e a dependência \omterm{def:b2:metric:continuity}{contínua} $\norm{y(t) -
z(t)} \leq \norm{y_0 - z_0}\,\eu^{k\abs{t - t_0}}$ para duas
soluções com dados iniciais diferentes.
\end{exercise}

\begin{exercise}[$\star\star$]\label{exo:b2:diffeq:7}
Sabendo que $y_1(t) = \frac{\sin t}{t}$ resolve $ty'' + 2y' + ty =
0$ em $\intoo{0}{\pi}$, ache uma segunda solução independente por
redução de \omterm{def:b2:structures:generated}{ordem} e dê a solução geral.
\end{exercise}

\begin{exercise}[$\star\star\star$]\label{exo:b2:diffeq:8}
Seja $A \in \mathcal{M}_n(\C)$ com todos os \omterm{def:b2:reduction:eigen}{autovalores} de parte real
(estritamente) negativa. Prove que toda solução de $X' = AX$ tende a
$0$ quando $t \to +\infty$, com taxa exponencial: $\norm{X(t)}
\leq C\,\eu^{-\alpha t}$ para algum $\alpha > 0$.
\emph{(Triangularize; trate o sistema triangular da última linha
para cima, ou use Dunford: $\eu^{tA} = \eu^{tD}\eu^{tN}$ com
$\norm{\eu^{tD}} \leq \eu^{-\alpha' t}$ e $\eu^{tN}$ polinomial
em $t$.)}
\end{exercise}

\begin{exercise}[$\star\star\star$]\label{exo:b2:diffeq:9}
(Sem fuga em tempo finito para crescimento linear) Suponha $f$
\omterm{def:b2:metric:continuity}{contínua} com $\norm{f(t, y)} \leq a\norm y + b$ em
$\intco{0}{\infty} \times \R^n$ e localmente \omterm{def:b2:metric:continuity}{lipschitziana} em $y$. Usando
Gronwall (\cref{exo:b2:diffeq:6}) na forma integral, prove que as
soluções maximais são globais (definidas em todo o
$\intco{0}{\infty}$).
\end{exercise}

\begin{exercise}[$\star$]\label{exo:b2:diffeq:10}
Resolva $y'' - 3y' + 2y = \eu^{t}$: soluções homogêneas e depois uma
solução particular da forma $\alpha t\,\eu^{t}$ \emph{(por que
o palpite ingênuo $\alpha\eu^t$ falha?)}; solução geral e
a solução com $y(0) = y'(0) = 0$.
\end{exercise}

\begin{exercise}[$\star\star$]\label{exo:b2:diffeq:11}
Calcule $\eu^{tA}$ para o bloco de Jordan
\[
A = \begin{pmatrix} \lambda & 1 & 0\\ 0 & \lambda & 1\\
0 & 0 & \lambda\end{pmatrix},
\]
e descreva todas as soluções de $X' = AX$: exponenciais vezes
vetores polinomiais, com graus até $2$. De onde vem o
grau polinomial?
\end{exercise}

\begin{exercise}[$\star\star\star$]\label{exo:b2:diffeq:12}
(Forçamento periódico, resposta periódica) Sejam $A \in
\mathcal{M}_n(\R)$ e $B \colon \R \to \R^n$ \omterm{def:b2:metric:continuity}{contínua} e
$T$-periódica.
\begin{enumerate}
  \item Mostre que uma solução de $X' = AX + B(t)$ é
        $T$-periódica se e somente se $X(T) = X(0)$
        \emph{(compare $X(\cdot + T)$ e $X$)}.
  \item Mostre que os \omterm{def:b2:reduction:eigen}{autovalores} de $\eu^{TA}$ são os
        $\eu^{T\lambda}$, $\lambda \in \operatorname{Sp}A$
        \emph{(triangularize sobre $\C$)}. Deduza: se nenhum
        \omterm{def:b2:reduction:eigen}{autovalor} de $A$ está em $\frac{2\iu\pi}{T}\Z$, então
        $I - \eu^{TA}$ é invertível.
  \item Sob essa hipótese, prove que o sistema tem
        exatamente uma solução $T$-periódica, com
        \[
        X(0) = \bigl(I - \eu^{TA}\bigr)^{-1}
        \int_0^{T}\eu^{(T-s)A}B(s)\,\dd s .
        \]
        A que corresponde o caso excluído, para o
        oscilador harmônico? (O problema de fim de semana responde:
        \omterm{pb:b2:diffeq:1}{ressonância}.)
\end{enumerate}
\end{exercise}

\section{Problema: oscilações, ressonância e os teoremas de
comparação de Sturm}

\begin{problem}\label{pb:b2:diffeq:1}
Uma equação governa o mundo mecânico e elétrico:
\[
x'' + 2\zeta\omega\,x' + \omega^2 x = F(t),
\qquad \omega > 0,\ \zeta \geq 0 .
\]
Este problema a estuda por completo --- pela
classificação traço--determinante dos sistemas lineares planos,
os três regimes de amortecimento, a resposta em regime permanente ao
forçamento periódico com seu pico de \emph{ressonância}\index{ressonância}
e a catástrofe de \omterm{pb:b2:diffeq:1}{ressonância} --- e depois deixa os coeficientes
constantes pelos \emph{teoremas de separação e comparação de
Sturm}\index{teorema de comparação de Sturm}, que controlam os
zeros das soluções de $y'' + q(t)y = 0$ sem fórmula
alguma.

\medskip
\noindent\textbf{Parte I --- O plano traço--determinante.} Sejam
$A \in \mathcal M_2(\R)$, $\tau = \operatorname{tr}A$, $\delta
= \det A$, $\Delta = \tau^2 - 4\delta$.
\begin{enumerate}
  \item Mostre que os \omterm{def:b2:reduction:eigen}{autovalores} de $A$ são
        $\frac{\tau\pm\sqrt\Delta}{2}$ e classifique: dois \omterm{def:b2:reduction:eigen}{autovalores}
        reais de sinais opostos se e somente se $\delta < 0$;
        \omterm{def:b2:reduction:eigen}{autovalores} reais de mesmo sinal se e somente se $\delta > 0$, $\Delta
        \geq 0$ (sinal de $\tau$); par conjugado não real se e somente se
        $\Delta < 0$ (parte real $\frac\tau2$).
  \item (Sela, $\delta < 0$) Com \omterm{def:b2:reduction:eigen}{autovalores} $\mu < 0 <
        \lambda$ e \omterm{def:b2:reduction:eigen}{autovetores} $v_\pm$, escreva a solução
        geral e descreva as trajetórias: duas semirretas estáveis e
        duas instáveis, e todas as demais órbitas assintóticas a
        ambas. Por que nenhuma solução além de $0$ pode permanecer limitada em todo
        o $\R$?
  \item (Nós, $\delta > 0$, $\Delta > 0$) Para $\mu < \lambda
        < 0$: mostre que toda solução não nula tende a $0$ e
        que todas as órbitas, exceto as do eixo rápido, chegam
        \emph{tangentes à direção própria lenta}
        \emph{(compare $\eu^{\mu t}$ e $\eu^{\lambda t}$)}.
  \item (Espirais e centros, $\Delta < 0$) Escrevendo os
        \omterm{def:b2:reduction:eigen}{autovalores} $\alpha \pm \iu\beta$, use
        \cref{exo:b2:diffeq:5} (após uma mudança de base real,
        admitida nessa generalidade ou demonstrada para os sistemas
        da Parte II, que são os usados adiante) para descrever
        as órbitas: espirais convergentes para $\alpha =
        \frac\tau2 < 0$, divergentes para $\tau > 0$, curvas
        fechadas (centro) para $\tau = 0$.
  \item (Casos de fronteira) Para um \omterm{def:b2:reduction:eigen}{autovalor} duplo ($\Delta =
        0$): mostre que $\eu^{tA} = \eu^{\lambda t}(I + tN)$ com $N
        = A - \lambda I$ nilpotente, e distinga a estrela
        ($N = 0$) do nó impróprio ($N \neq 0$).
        Resuma a Parte I no retrato traço--determinante da
        figura deste capítulo.
\end{enumerate}

\medskip
\noindent\textbf{Parte II --- O oscilador amortecido.} Agora $F =
0$: $x'' + 2\zeta\omega x' + \omega^2x = 0$, isto é, $X' = AX$
com $A = \begin{pmatrix} 0 & 1\\ -\omega^2 &
-2\zeta\omega\end{pmatrix}$.
\begin{enumerate}[resume]
  \item Calcule $\tau, \delta, \Delta$ e situe os três
        regimes no plano traço--determinante:
        \emph{subamortecido} $0 < \zeta < 1$ (espiral estável),
        \emph{criticamente amortecido} $\zeta = 1$ (\omterm{def:b2:reduction:eigen}{autovalor}
        duplo), \emph{superamortecido} $\zeta > 1$ (nó
        estável); $\zeta = 0$ é o centro.
  \item Resolva os três regimes explicitamente:
        \[
        \zeta < 1:\ \eu^{-\zeta\omega t}\bigl(a\cos\omega_d t
        + b\sin\omega_dt\bigr),\ \omega_d =
        \omega\sqrt{1-\zeta^2};
        \qquad
        \zeta = 1:\ (a + bt)\,\eu^{-\omega t};
        \]
        $\zeta > 1$: duas exponenciais reais. Defina o
        pseudoperíodo $\frac{2\pi}{\omega_d}$ e mostre que
        a razão entre máximos sucessivos de $\abs x$ é a
        constante $\eu^{-2\pi\zeta/\sqrt{1-\zeta^2}}$ (o
        decremento logarítmico).
  \item (O princípio da mola de porta) Para $\zeta \geq 1$ a
        taxa de decaimento é governada pelo \omterm{def:b2:reduction:eigen}{autovalor} mais lento
        $\lambda_{\mathrm{slow}} = -\omega\bigl(\zeta -
        \sqrt{\zeta^2-1}\bigr)$. Mostre que
        $\abs{\lambda_{\mathrm{slow}}} =
        \frac{\omega}{\zeta + \sqrt{\zeta^2 - 1}}$ é função
        \emph{decrescente} de $\zeta \geq 1$:
        o amortecimento crítico $\zeta = 1$ dá o retorno ao
        repouso mais rápido sem oscilação.
  \item (Energia) Seja $E(t) = \frac12x'^2 +
        \frac12\omega^2x^2$. Prove que $E' = -2\zeta\omega\,x'^2
        \leq 0$ e deduza que, para $\zeta > 0$, a equação
        não tem solução periódica não nula \emph{(um período
        forçaria $E$ constante, logo $x' \equiv 0$)}.
  \item Explique em duas frases por que o centro $\zeta = 0$
        é \emph{estruturalmente frágil}: qualquer $\zeta > 0$,
        por menor que seja, destrói a periodicidade --- e onde
        isso aparece no plano traço--determinante (a reta dos
        centros tem interior vazio).
\end{enumerate}

\medskip
\noindent\textbf{Parte III --- Oscilações forçadas e
\omterm{pb:b2:diffeq:1}{ressonância}.} Agora $F(t) = F\cos(\gamma t)$ com $F, \gamma > 0$.
\begin{enumerate}[resume]
  \item ($\zeta > 0$: o regime permanente) Procure $x_p =
        \Re\bigl(z\,\eu^{\iu\gamma t}\bigr)$: mostre que
        \[
        z = \frac{F}{\omega^2 - \gamma^2 +
        2\iu\zeta\omega\gamma},
        \qquad
        A(\gamma) := \abs z = \frac{F}{\sqrt{(\omega^2 -
        \gamma^2)^2 + 4\zeta^2\omega^2\gamma^2}} ,
        \]
        e escreva $x_p = A(\gamma)\cos(\gamma t - \varphi)$
        com $\tan\varphi =
        \frac{2\zeta\omega\gamma}{\omega^2-\gamma^2}$.
  \item Mostre que \emph{toda} solução é $x_p$ mais um
        transiente da Parte II, que tende a $0$: quaisquer que
        sejam os dados iniciais, o sistema se trava no regime
        permanente --- amplitude $A(\gamma)$, atraso de fase
        $\varphi$.
  \item (A curva de \omterm{pb:b2:diffeq:1}{ressonância}) Maximize $A$: mostre que
        $A(\gamma)$ tem máximo interior se e somente se $\zeta <
        \frac{1}{\sqrt2}$, em
        \[
        \gamma_* = \omega\sqrt{1 - 2\zeta^2},
        \qquad
        A(\gamma_*) =
        \frac{F}{2\zeta\omega^2\sqrt{1-\zeta^2}} ,
        \]
        e que, para $\zeta$ pequeno, o pico amplifica a
        resposta estática $A(0) = \frac F{\omega^2}$ pelo
        fator $\approx \frac{1}{2\zeta}$.
  \item ($\zeta = 0$, fora da \omterm{pb:b2:diffeq:1}{ressonância}) Para $\gamma \neq \omega$,
        mostre que a solução com $x(0) = x'(0) = 0$ é
        \[
        x(t) = \frac{F}{\omega^2 -
        \gamma^2}\bigl(\cos\gamma t - \cos\omega t\bigr)
        = \frac{2F}{\omega^2-\gamma^2}
        \sin\frac{(\omega-\gamma)t}{2}
        \sin\frac{(\omega+\gamma)t}{2} :
        \]
        limitada, com \emph{batimentos} --- uma oscilação rápida
        sob um envelope lento --- quando $\gamma$ está próximo de
        $\omega$.
  \item ($\zeta = 0$, \omterm{pb:b2:diffeq:1}{ressonância}) Para $\gamma = \omega$, mostre
        que $x_p(t) = \frac{F}{2\omega}\,t\sin(\omega t)$ é
        solução e recupere-a como limite da questão
        14 quando $\gamma \to \omega$: a amplitude cresce
        linearmente para sempre --- a catástrofe de \omterm{pb:b2:diffeq:1}{ressonância}.
  \item (Elo com Fourier) Um forçamento periódico geral se decompõe
        em harmônicos (o capítulo de Fourier); por linearidade, o
        regime permanente é a soma das respostas harmônicas. Para
        um oscilador não amortecido de frequência $\omega = 3$
        forçado pela onda quadrada do tipo
        \cref{exo:b2:fourier:1} (harmônicos em todos os
        inteiros ímpares), qual harmônico entra em \omterm{pb:b2:diffeq:1}{ressonância}? Uma frase sobre
        por que os engenheiros temem ondas quadradas.
\end{enumerate}

\medskip
\noindent\textbf{Parte IV --- Os teoremas de Sturm.} Considere $y''
+ q(t)\,y = 0$ num intervalo $I$, com $q$ \omterm{def:b2:metric:continuity}{contínua}. (Toda
equação $y'' + ay' + by = 0$ reduz-se a essa forma normal pela
substituição $y = u\exp\bigl(-\frac12\int a\bigr)$;
a questão 21 mostra uma variante do truque em ação.)
\begin{enumerate}[resume]
  \item Para duas soluções $y_1, y_2$, mostre que o \omterm{def:b2:diffeq:wronskian}{wronskiano}
        $W = y_1y_2' - y_1'y_2$ é \emph{constante}, nulo se e somente se
        as soluções são proporcionais; e que uma solução não nula
        tem apenas zeros \emph{simples e isolados}.
  \item (Separação de Sturm) Sejam $y_1, y_2$ soluções independentes
        e $a < b$ dois zeros consecutivos de $y_1$.
        Prove que $y_2$ se anula exatamente uma vez em
        $\intoo{a}{b}$ \emph{(avalie a constante $W$ em $a$
        e $b$: $W = y_1'y_2$ aí, e $y_1'(a)$,
        $y_1'(b)$ têm sinais opostos)}: os zeros de soluções
        independentes se intercalam.
  \item (Comparação de Sturm) Sejam $q_1 \leq q_2$ em $I$, $y \neq
        0$ com $y'' + q_1y = 0$, $z \neq 0$ com $z'' + q_2z
        = 0$, e $a < b$ zeros consecutivos de $y$. Mostre que
        $z$ se anula em $\intcc{a}{b}$ --- estritamente dentro se
        $q_1 < q_2$ em algum ponto de $\intoo ab$ \emph{(se $z \neq
        0$ em $\intoo ab$, estude $(yz' - y'z)' = (q_1 -
        q_2)yz$ com sinais fixos para $y, z$ e compare os
        valores de bordo)}.
  \item Deduza as \emph{cotas de espaçamento}: se $0 < m^2 \leq
        q(t) \leq M^2$ em $I$, então dois zeros consecutivos
        $a < b$ de uma solução não nula de $y'' + qy = 0$
        satisfazem
        \[
        \frac{\pi}{M} \;\leq\; b - a \;\leq\; \frac{\pi}{m}
        \]
        \emph{(compare com $u'' + M^2u = 0$ e $u'' + m^2u
        = 0$, cujos zeros distam $\frac\pi M$ e
        $\frac\pi m$)}. Confira no oscilador harmônico.
  \item Transforme $ty'' + 2y' + ty = 0$
        (\cref{exo:b2:diffeq:7}) por $u = ty$ em $u'' + u =
        0$, recupere instantaneamente suas soluções $\frac{\sin t}t$,
        $\frac{\cos t}{t}$ e conclua que os
        zeros de toda solução não nula distam exatamente
        $\pi$: a visão de mundo de Sturm --- os zeros são controlados pelo
        coeficiente $q$, com fórmula ou sem.
\end{enumerate}

\medskip
\noindent\textbf{Parte V --- Duhamel e a fronteira da
limitação.}
\begin{enumerate}[resume]
  \item (Duhamel para o oscilador) Mostre que, para $F$
        \omterm{def:b2:metric:continuity}{contínua}, a solução de $x'' + \omega^2x =
        F(t)$ com $x(0) = x'(0) = 0$ é
        \[
        x(t) = \frac1\omega\int_0^t\sin\bigl(\omega(t -
        s)\bigr)F(s)\,\dd s ,
        \]
        e rededuza dela a solução ressonante da questão 15
        com $F(s) = F\cos(\omega s)$
        \emph{(produto-soma)}.
  \item ($\zeta > 0$: entrada limitada, saída limitada) Mostre
        que, para $\zeta > 0$ e \emph{qualquer} $F$ \omterm{def:b2:metric:continuity}{contínua}
        limitada, toda solução da equação amortecida é limitada
        em $\intco{0}{\infty}$ \emph{(variação das constantes
        mais o decaimento exponencial $\vertiii{\eu^{tA}} \leq
        C\eu^{-\alpha t}$ de \cref{exo:b2:diffeq:8})}.
  \item ($\zeta = 0$) Mostre que, sem amortecimento, um forçamento
        periódico limitado mantém todas as soluções limitadas
        \emph{exceto} exatamente na \omterm{pb:b2:diffeq:1}{ressonância} ($\gamma =
        \omega$, questão 15 contra questão 14): o amortecimento é
        o que transforma a fronteira da limitação em estabilidade
        uniforme.
  \item Síntese. Uma frase para cada: (i) como o
        plano traço--determinante organiza as Partes I--II e
        onde o forçamento (Parte III) o deixa; (ii) o
        significado físico de $\gamma_*$, $A(\gamma_*)$ e do
        fator $\frac1{2\zeta}$; (iii) o que os teoremas de Sturm
        dizem que as fórmulas explícitas não conseguem dizer; (iv)
        quais dois resultados deste problema o resto do
        livro vai reutilizar em silêncio (os \omterm{def:b2:diffeq:wronskian}{wronskianos} de
        constante de Liouville; a estabilidade com entrada limitada).
\end{enumerate}
\end{problem}
